\chapter{偶像}
\section{偶像的特征}

\begin{itemize}
\item 偶像不仅仅是有形的，也包括无形的\\
\begin{itemize}
\item 以弗所書 5 因为你们确实地知道，无论是淫乱的，是污秽的，是有贪心的，在基督和 神的国里都是无份的。\textcolor{red}{有贪心的，就与拜偶像的一样}。
\item 歌羅西書 3:5 所以要治死你們在地上的肢體，就如淫亂、汙穢、邪情、惡慾和貪婪，貪婪就與拜偶像一樣。
\item 提摩太前書 6:10,17 貪財是萬惡之根。有人貪戀錢財，就被引誘離了真道，用許多愁苦把自己刺透了。
\item 馬太福音 6:24 一個人不能侍奉兩個主，不是惡這個愛那個，就是重這個輕那個。你們不能又侍奉神，又侍奉瑪門。
\end{itemize}
\item 偶像什么也不能做\\
詩篇115:3-7 世人的偶像是金的銀的，是人用手做的。
偶像有嘴巴，卻不能說話；
有眼睛，卻不能看見；
有耳朵，卻聽不見；
有鼻子，卻聞不到；
有手，卻不能觸摸；
有腳，卻不能走路；
有喉嚨，也不能發聲。
\item 偶像不能降禍，也無力降福\\
耶利米書 10:3-5 眾民的風俗是虛空的，他們在樹林中用斧子砍伐一棵樹，匠人用手工造成偶像。 4 他們用金銀裝飾它，用釘子和錘子釘穩，使它不動搖。 5 它好像棕樹，是鏇成的，不能說話不能行走，必須有人抬著。你們不要怕它，它不能降禍，也無力降福。」
\end{itemize}

\section{拜偶像人的特征——被偶像迷住不能自拔，有心不能思想(詩篇115:8)}
製造偶像的人，
會跟偶像一樣；
信賴偶像的，
都會與偶像相同。


\section{向拜偶像人传福音}
\subsection{保罗向拜偶像人传福音的策略 (使徒行传17:16-34)}


\subsubsection{保罗在雅典发现到处都是偶像 (使徒行传17:16-34)}
\begin{table}[H]
\centering
\begin{tabular}{p{17.25cm}}\hline\hline
16 保羅在雅典等他們（提摩太，路加，希拉）的時候，看見城裡到處都是偶像，就十分不安。17 於是他\textcolor{red}{在會堂裡跟猶太人和其他崇拜上帝的人推理}，\\\hline
也天天\textcolor{red}{在市集廣場上跟遇見的人推理}。18 可是，有些伊壁鳩魯派和斯多葛派的哲學家來跟他辯論。有的說：「這個人囉囉唆唆，到底想說什麼？」有的說：「看來他是在傳揚外國的神明。」他們這樣說，是因為保羅宣揚耶穌和復活的好消息。\\\hline
19 他們把保羅帶到\textcolor{red}{亞略巴古}，說：「你講的這套新道理，能說給我們聽聽嗎？20 你講的這些事，聽起來很新奇，我們想知道是什麼意思。」21 原來所有雅典人和暫住當地的外國人，有空時最喜歡聊聊天，聽聽新鮮事。\\\hline
\end{tabular}
\end{table}
\begin{figure}[H]
\centering
\includegraphics[width=6.in]{ZongJie/figs/YiXiJie/OuXiang1}
\end{figure}

\subsubsection{保罗的讲道 (使徒行传17:22-34)}
\begin{itemize}
\item 正面肯定雅典人寻求神的优点\\
22 保羅站在亞略巴古中間，說：「各位雅典人，我看出\textcolor{red}{你們事事都似乎比其他人更敬畏神明}。23 比如說，我一邊走一邊觀察你們崇敬的事物，發現有一座壇，上面寫著『獻給不認識的神』。
\item 传讲保罗所信的神的特征\\
現在我要向你們傳揚的，其實就是你們不認識卻崇拜的神。24 他是創造世界和世上萬物的真神，既然是天地的主，當然不會住在人建造的殿宇裡，25 也不需要人服侍，好像缺什麼似的，因為把生命、氣息和一切賜給所有人的是他。26 他從一個人造出了萬國萬族，讓他們住在大地上，他定下時期和人們居住的疆界，27 讓他們可以尋求真神，只要努力尋找就能找到，其實他離我們每個人都不遠。
\item 从雅典诗人的诗句和真神的关系入手，劝导众人悔改信真神\\
28 我們能有生命，能活動，能存在，都有賴於他，就像\textcolor{red}{你們有些詩人說的一樣：『我們也是他的孩子}。』 29 「我們\underline{既然是真神的孩子}，就不該以為他好像金子、銀子或石頭，好像人用手藝和心思雕塑出來的東西。30 在以往的時代，世人這麼無知，真神並沒有追究。現在他卻向世界各地的人宣告，他們必須悔改，31 因為他已經定下日子，決定通過他委任的人審判全天下，伸張正義。他已經使這個人死而復生，以此作為對所有人的保證。」
\end{itemize}

\subsubsection{讲道的结果，有些人信主 (使徒行传17:32-34)}
32 有些人聽到死人復活這件事，就開始嘲笑他，有些人卻說：「我們下次要再聽你講講這件事。」33 於是保羅離開了他們。34 有些人卻跟他來往，成了信徒，其中包括亞略巴古法庭的審判官丟尼修，還有一個名叫達馬里的女人以及其他人。

\subsection{保罗向拜偶像人传福音的策略给我们的启发}
使徒行传记载了保罗在希腊的雅典，向拜偶像的人成功传福音的策略，给我的启发如下供你参考。
\begin{itemize}
\item 正面肯定拜偶像之人有一颗想寻求神的优点，如保罗对雅典人所做的。因为这样做，可以取消对方心中对我们的敌意，挪去偶像在人心中所设的拦阻。对方容易安静心听我们要讲的。
\item 愿主赐给你智慧，传讲我们所信的神的特征。保罗的讲论不长，但句句射中靶心。他是针对雅典人信的假神的特征，传讲真神的特征。不同的人群，被捆绑的方式不一样。
\item 从对方熟悉的文化或者知识背景入手，找出和我们信的神的关系，带出劝人悔改的信息。
\end{itemize}





